\section{Conclusion}
\label{sec:conclusion}

We introduced SEAL, a method that groups longitudinal entities by the shape of their target trajectory and describes each group by a short rule, treating the grouping and the description as equal partners joined by one cost. In keeping with the descriptive-mining tradition, we assessed it by whether its rules recover structure that domain readers already accept, and across the four public panels they did. Per-capita emission trajectories separated by national income and fuel mix, placing the Gulf producers and North America in one group and the coal-reliant industrial economies in another. Cross-country mortality trajectories separated first by the age structure of the population and then, by nothing more than population size and density, into recognisable regional waves. Income trajectories separated the Asian catch-up economies from the resource-dependent boom-and-bust economies by the shape of their growth rather than its endpoint. Particulate trajectories separated by geography even though land-use labels were available to the search. This agreement is the main evidence that the rules are meaningful rather than accidental, since the drivers were never supplied to the method.

Two aspects of the preparation carried much of the weight. Turning time-varying covariates into engineered temporal descriptors let the rules speak about how a covariate moved, for example a birth rate that was low and still falling, rather than only about its ending value. Holding the descriptors to a non-circular discipline kept the rules from restating the definition of their own target, which was essential in the household health panel where the target is a ratio of two quantities that could otherwise leak into the descriptors.

The method has clear limits. The distance between trajectories is computed for every pair of entities, so the cost grows with the square of the panel size, which kept the candidate budget modest and would need a different distance or an approximation to reach much larger panels. The cumulative representation emphasises persistent level differences and is less sensitive to short reversals, which suits some targets better than others. The household survey is not a population sample, so its findings describe the surveyed households rather than the country.

Several directions follow naturally. The household health panel invites a second stage that separates subgroups whose burden moved because income moved, which is mechanical through the denominator, from subgroups whose burden moved because spending behaviour changed at steady income, which is the more actionable finding. Scaling beyond the quadratic distance, extending the temporal descriptors, and handling discrete or event-like targets would each widen the range of panels the method can serve. Our results suggest that searching for a grouping and its rule together, rather than in separate steps, is a practical way to read structure out of longitudinal data in terms a domain audience can use.
